% Generated draft: requires mathematical and source review.
\documentclass[11pt]{article}
\usepackage[margin=25mm]{geometry}
\usepackage{fontspec}
\setmainfont{Latin Modern Roman}
\usepackage{amsmath,amssymb,mathtools,unicode-math}
\setmathfont{Latin Modern Math}
\usepackage{enumitem,needspace,xurl,hyperref}
\setlength{\parindent}{0pt}
\setlength{\parskip}{5pt}
\setlength{\emergencystretch}{3em}
\newcommand{\N}{\mathbb{N}}
\newcommand{\Z}{\mathbb{Z}}
\newcommand{\Q}{\mathbb{Q}}
\newcommand{\R}{\mathbb{R}}
\newcommand{\C}{\mathbb{C}}
\newcommand{\F}{\mathbb{F}}
\newcommand{\A}{\mathbb{A}}
\newcommand{\PP}{\mathbb{P}}
\newcommand{\E}{\mathbb{E}}
\newcommand{\eps}{\varepsilon}
\newcommand{\dd}{\,\mathrm{d}}
\begin{document}
\section*{UnsolvedMath ID: 701427}
\textbf{Source problem code:} catalogue entry
\subsection{Definitions and mathematical statement}
For elements of a group, write \([x,y]=x^{-1}y^{-1}xy\), and let \([G,G]=\langle [x,y]:x,y\in G\rangle\) be the derived subgroup. A group \(G\) is residually finite if, for each \(g\ne 1\), there is a homomorphism \(\varphi:G\to F\) to a finite group with \(\varphi(g)\ne 1\). A group is locally finite if every finitely generated subgroup is finite. The identity \([x,y]^n=1\) means that this equation holds for all \(x,y\in G\). We use the usual interpretation that \(n\) is a positive integer; the supplied statement does not explicitly specify its range.

\textbf{Main question.} For every positive integer \(n\) and every residually finite group \(G\) satisfying \([x,y]^n=1\) for all \(x,y\in G\), must \([G,G]\) be locally finite?

A group is soluble if its derived series, defined by \(G^{(0)}=G\) and \(G^{(i+1)}=[G^{(i)},G^{(i)}]\), reaches the trivial subgroup. For a finite soluble group \(G\), its Fitting height \(h(G)\) is the least integer \(h\geq 0\) for which there is a series
\[
1=G_0\leq G_1\leq\cdots\leq G_h=G,
\]
with every \(G_i\) normal in \(G\) and every factor \(G_i/G_{i-1}\) nilpotent. A group is nilpotent if its lower central series \(\gamma_1(H)=H\), \(\gamma_{i+1}(H)=[\gamma_i(H),H]\) reaches the trivial subgroup, where \([A,B]=\langle[a,b]:a\in A,\ b\in B\rangle\). In particular, \(h(1)=0\).

\textbf{Alternative question, described as equivalent in the supplied source.} For every positive integer \(n\), does there exist an integer \(H_n\geq 0\) such that every finite soluble group \(G\) satisfying \([x,y]^n=1\) for all \(x,y\in G\) has \(h(G)\leq H_n\)? Thus the bound may depend on \(n\), but not on \(G\). The equivalence is a supplied source assertion; no justification is included in the record. The source also says ``Cf. Archive, 13.34'', but the referenced archive entry is not supplied.
\subsection{Short English statement}
If all commutators in a residually finite group have orders dividing a fixed positive integer \(n\), must every finitely generated subgroup of its derived subgroup be finite? The supplied source describes this as equivalent to asking for a bound, depending only on \(n\), on the Fitting heights of finite soluble groups satisfying the same identity.
\subsection{Sources}
\begin{itemize}
\item Local UnsolvedMath registry, permanent ID \(701427\), exact source problem code \texttt{catalogue entry}, title \emph{Kourovka Notebook Problem 17.126}. The supplied registry labels the problem \texttt{open}; this label has not been independently verified.
\item \emph{The Kourovka Notebook (Unsolved Problems in Group Theory)}, problem 17.126, proposed by P. Shumyatsky in 2010. The supplied background identifies the 17th issue (2010) and the 21st issue edition (2026), arXiv identifier \texttt{1401.0300v48}. The supplied collection description names E. I. Khukhro and V. D. Mazurov as editors.
\item Supplied cross-reference: ``Cf. Archive, 13.34''. The referenced entry and its relationship to this problem are not included in the supplied record.
\end{itemize}
\end{document}
